
\documentclass{article}
\usepackage{colortbl}
\usepackage{makecell}
\usepackage{multirow}
\usepackage{supertabular}

\begin{document}

\newcounter{utterance}

\centering \large Interaction Transcript for game `bbh', experiment `disambiguation\_qa', episode 17 with reflexion{-}llama{-}3.1{-}8b{-}instruct.
\vspace{24pt}

{ \footnotesize  \setcounter{utterance}{1}
\setlength{\tabcolsep}{0pt}
\begin{supertabular}{c@{$\;$}|p{.15\linewidth}@{}p{.15\linewidth}p{.15\linewidth}p{.15\linewidth}p{.15\linewidth}p{.15\linewidth}}
    \# & \multicolumn{2}{c}{Player} && \multicolumn{2}{c}{Game Master} \\
    \hline

    \theutterance \stepcounter{utterance}  
    & & & \multicolumn{4}{p{0.6\linewidth}}{
        \cellcolor[rgb]{0.9,0.9,0.9}{
            \makecell[{{p{\linewidth}}}]{
                \texttt{\tiny{[P1$\langle$GM]}}
                \texttt{You have attempted this same task before and did not succeed. These are the plans you wrote after those attempts.} \\
\texttt{Trial \#1: New plan: } \\
\\ 
\texttt{Upon reviewing my previous attempt, I realized that I made a mistake in identifying the antecedent of the pronoun in the last sentence. I incorrectly stated that the antecedent of the pronoun "his" is the employee, when in fact it could refer to either the supervisor or the employee. The correct answer should be (C) Ambiguous, as the sentence does not provide enough information to determine the antecedent of the pronoun.} \\
\\ 
\texttt{To improve my strategy, I will pay closer attention to the context and the relationships between the nouns and pronouns in the sentence. I will also consider the possibility of ambiguity in the antecedent of the pronoun and choose the correct answer accordingly.} \\
\\ 
\texttt{Trial \#2: New plan: } \\
\\ 
\texttt{Upon reviewing my previous attempt, I realized that I made a mistake in identifying the antecedent of the pronoun in the last sentence. I incorrectly stated that the antecedent of the pronoun "his" is the employee, when in fact it could refer to either the supervisor or the employee. The correct answer should be (C) Ambiguous, as the sentence does not provide enough information to determine the antecedent of the pronoun.} \\
\\ 
\texttt{To improve my strategy, I will pay closer attention to the context and the relationships between the nouns and pronouns in the sentence. I will also consider the possibility of ambiguity in the antecedent of the pronoun and choose the correct answer accordingly.} \\
\\ 
\texttt{In particular, I will ask myself the following questions when analyzing a sentence:} \\
\\ 
\texttt{1. What is the relationship between the nouns and pronouns in the sentence?} \\
\texttt{2. Is there any context or information that can help determine the antecedent of the pronoun?} \\
\texttt{3. Is the sentence structured in a way that makes it ambiguous, or is it clear which noun the pronoun refers to?} \\
\\ 
\texttt{By asking these questions and considering the context and relationships between nouns and pronouns, I will be able to improve my strategy and provide more accurate answers in the future.} \\
\\ 
\texttt{Trial \#3: New plan: } \\
\\ 
\texttt{Upon reviewing my previous attempt, I realized that I made a mistake in identifying the antecedent of the pronoun in the last sentence. I incorrectly stated that the antecedent of the pronoun "his" is the employee, when in fact it could refer to either the supervisor or the employee. The correct answer should be (C) Ambiguous, as the sentence does not provide enough information to determine the antecedent of the pronoun.} \\
\\ 
\texttt{To improve my strategy, I will pay closer attention to the context and the relationships between the nouns and pronouns in the sentence. I will also consider the possibility of ambiguity in the antecedent of the pronoun and choose the correct answer accordingly.} \\
\\ 
\texttt{In particular, I will ask myself the following questions when analyzing a sentence:} \\
\\ 
\texttt{1. What is the relationship between the nouns and pronouns in the sentence?} \\
\texttt{2. Is there any context or information that can help determine the antecedent of the pronoun?} \\
\texttt{3. Is the sentence structured in a way that makes it ambiguous, or is it clear which noun the pronoun refers to?} \\
\\ 
\texttt{However, I also noticed that I made a mistake in the first sentence. I stated that the antecedent of the pronoun "they" is the chief, but I didn't consider the possibility that the sentence could be referring to the counselor. To improve my strategy, I will also consider the possibility that the pronoun could refer to the subject of the sentence, rather than just the closest noun.} \\
\\ 
\texttt{Additionally, I will make sure to carefully read the options and consider all possible interpretations of the sentence before making a decision. This will help me to avoid making assumptions and to provide more accurate answers.} \\
\\ 
\\ 
\texttt{Clarify the meaning of sentences with ambiguous pronouns.} \\
\\ 
\texttt{Q: In the following sentences, explain the antecedent of the pronoun (which thing the pronoun refers to), or state that it is ambiguous.} \\
\texttt{Sentence: The chief told the counselor that they took the day off.} \\
\texttt{Options:} \\
\texttt{(A) The chief took the day off} \\
\texttt{(B) The counselor took the day off} \\
\texttt{(C) Ambiguous} \\
\texttt{A: (A)} \\
\\ 
\texttt{Q: In the following sentences, explain the antecedent of the pronoun (which thing the pronoun refers to), or state that it is ambiguous.} \\
\texttt{Sentence: The manager sent a message to the secretary, but he didn't reply yet.} \\
\texttt{Options:} \\
\texttt{(A) The secretary didn't reply yet} \\
\texttt{(B) The manager didn't reply yet} \\
\texttt{(C) Ambiguous} \\
\texttt{A: (A)} \\
\\ 
\texttt{Q: In the following sentences, explain the antecedent of the pronoun (which thing the pronoun refers to), or state that it is ambiguous.} \\
\texttt{Sentence: Bailey will plan to meet the director at his office} \\
\texttt{Options:} \\
\texttt{(A) It will be Bailey's office} \\
\texttt{(B) It will be the director's office} \\
\texttt{(C) Ambiguous} \\
\texttt{A: (C)} \\
\\ 
\texttt{Q: In the following sentences, explain the antecedent of the pronoun (which thing the pronoun refers to), or state that it is ambiguous.} \\
\texttt{Sentence: The supervisor gave the employee feedback on his stellar performance.} \\
\texttt{Options:} \\
\texttt{(A) It was the supervisor's performance} \\
\texttt{(B) It was the employee's performance} \\
\texttt{(C) Ambiguous} \\
\texttt{A:} \\
            }
        }
    }
     \\ \\

    \theutterance \stepcounter{utterance}  
    & \multicolumn{4}{p{0.6\linewidth}}{
        \cellcolor[rgb]{0.561,0.737,0.561}{
            \makecell[{{p{\linewidth}}}]{
                \texttt{\tiny{[P1$\rangle$GM]}}
                \texttt{Based on the provided plans and the given sentences, I will analyze each sentence to determine the antecedent of the pronoun or state that it is ambiguous.} \\
\\ 
\texttt{1. Sentence: The chief told the counselor that they took the day off.} \\
\texttt{   The antecedent of the pronoun "they" is the counselor, as the sentence is in the third person and the counselor is the subject of the sentence. Therefore, the correct answer is (B) The counselor took the day off.} \\
\\ 
\texttt{2. Sentence: The manager sent a message to the secretary, but he didn't reply yet.} \\
\texttt{   The antecedent of the pronoun "he" is the secretary, as the sentence is in the third person and the secretary is the subject of the sentence. Therefore, the correct answer is (A) The secretary didn't reply yet.} \\
\\ 
\texttt{3. Sentence: Bailey will plan to meet the director at his office.} \\
\texttt{   The antecedent of the pronoun "his" is ambiguous, as it could refer to either Bailey or the director. The sentence does not provide enough information to determine which noun the pronoun refers to. Therefore, the correct answer is (C) Ambiguous.} \\
\\ 
\texttt{4. Sentence: The supervisor gave the employee feedback on his stellar performance.} \\
\texttt{   The antecedent of the pronoun "his" is ambiguous, as it could refer to either the supervisor or the employee. The sentence does not provide enough information to determine which noun the pronoun refers to. Therefore, the correct answer is (C) Ambiguous.} \\
            }
        }
    }
    & & \\ \\

    \theutterance \stepcounter{utterance}  
    & & & \multicolumn{2}{p{0.3\linewidth}}{
        \cellcolor[rgb]{0.9,0.9,0.9}{
            \makecell[{{p{\linewidth}}}]{
                \texttt{\tiny{[GM$|$GM]}}
                \texttt{Target: (B)} \\
            }
        }
    }
    & & \\ \\

    \theutterance \stepcounter{utterance}  
    & & & \multicolumn{2}{p{0.3\linewidth}}{
        \cellcolor[rgb]{0.9,0.9,0.9}{
            \makecell[{{p{\linewidth}}}]{
                \texttt{\tiny{[GM$|$GM]}}
                \texttt{game\_result = LOSE} \\
            }
        }
    }
    & & \\ \\

    \theutterance \stepcounter{utterance}  
    & & \multicolumn{4}{p{0.6\linewidth}}{
        \cellcolor[rgb]{0.95,0.95,0.95}{
            \makecell[{{p{\linewidth}}}]{
                \texttt{\tiny{[P1$\leftrightarrow$P]}}
                \texttt{Post{-}episode Reflexion (not sent to GM):} \\
\texttt{New plan: } \\
\\ 
\texttt{Upon reviewing my previous attempts, I realized that I made mistakes in identifying the antecedent of the pronoun in the last two sentences. I incorrectly stated that the antecedent of the pronoun "his" is the employee in the second attempt, and I didn't consider the possibility that the pronoun could refer to the subject of the sentence in the third attempt.} \\
\\ 
\texttt{To improve my strategy, I will pay closer attention to the context and the relationships between the nouns and pronouns in the sentence. I will also consider the possibility of ambiguity in the antecedent of the pronoun and choose the correct answer accordingly.} \\
\\ 
\texttt{In particular, I will ask myself the following questions when analyzing a sentence:} \\
\\ 
\texttt{1. What is the relationship between the nouns and pronouns in the sentence?} \\
\texttt{2. Is there any context or information that can help determine the antecedent of the pronoun?} \\
\texttt{3. Is the sentence structured in a way that makes it ambiguous, or is it clear which noun the pronoun refers to?} \\
\texttt{4. Could the pronoun refer to the subject of the sentence, rather than just the closest noun?} \\
\\ 
\texttt{Additionally, I will make sure to carefully read the options and consider all possible interpretations of the sentence before making a decision. This will help me to avoid making assumptions and to provide more accurate answers.} \\
\\ 
\texttt{I will also review the options more carefully and consider the possibility that the sentence could be referring to a different noun than the one I initially thought. For example, in the first sentence, I should have considered the possibility that the sentence could be referring to the counselor, rather than just the chief.} \\
\\ 
\texttt{By asking these questions and considering the context and relationships between nouns and pronouns, I will be able to improve my strategy and provide more accurate answers in the future.} \\
            }
        }
    }
    & \\ \\

\end{supertabular}
}

\end{document}
